\Titular*%
{profesorado}
{O meu rexistro de frases de profesores}%
{Manuel Galán Rodríguez}%
{Catro anos apuntando as xoias da facultade.}

\noindent\begin{center}\fbox{
    \parbox{0.85\linewidth}{
        Á vista de non dispor dun artigo dalgún dos nosos profesores,
        escollemos que aquí o presente redactor plasmase algunhas frases
        memorábeis dos mesmos.
    }
}\end{center}
\vspace*{5mm}
\begin{multicols}{2}
\raggedcolumns

Para empezar, supoño que ti, querido lector, terás ouvido falar algunha vez do
meu rexistro de frases célebres de profesores da Facultade de Física da USC que
teño na miña conta persoal de X (\texttt{@manolo\_o\_galan}), que animo que
leas, en caso de que estiveses interesado. En caso afirmativo, entendo que
tamén che chamará a atención este artigo, co que te convido a que sigas lendo.

Unha primeira pregunta que che pode estar xurdindo neste momento é, como se me
deu por levar a cabo todo este rexistro? Se cadra, algún dos máis veteranos
estará pensando que a recompilación de frases que realizou David Castaño, xa
hai uns cantos anos, puido ter algo que ver, aínda que a realidade é bastante
máis complexa como detallarei a continuación.

No meu percorrido polo Sistema Educativo Español, á altura do curso de 4ºESO,
coincidiume de ter varios profesores bastante simpáticos, cuxas clases eran
bastante amenas. Un día, ocorréuseme que quería manter un rexistro para que así
cando fose maior puidese rememorar vellos tempos. Entre estas xoias podemos
destacar \textit{``Manuel, eu nunca reflexionei tanto sobre a acentuación!''}
da profesora de Galego ou a de \textit{``Puerta OR, como la matrícula de los
coches de Orense''} do de Tecnoloxía, falando de portas lóxicas.

Esta recompilación evidentemente non foi levada a cabo mediante a plataforma
Twitter, pois eu aínda era un rapaz de 15 anos, demasiado novo para esas
cousas.
%\footnote{Se tendes algún irmán ou primo con esa idade, por favor, que non se
%meta nesa clase de redes sociais, en Twitter está o peor do ser humano.}.
O lugar escollido foi un arquivo en Microsoft Word que tiña no meu computador e
que aínda se encontra polas súas profundidades. De cando en vez, cando pasa
pola miña cabeza ese asunto, vou e consúltoo, e aínda paso un bo tempo rindo.

Esta colección de xoias mantívose despois en 1ºBAC, aínda que con menor
intensidade. Despois, xa veu o confinamento, que o seguiu un 2ºBAC moito máis
aborrecido, que sentou mal tanto a profesores como a alumnos, o que supuxo unha
interrupción de devandita recompilación.

Cando estaba, non lembro ben se a finais de 2ºBAC estudando Historia de España,
ou se xa a comezos do 1º curso do Grao en Física, apareceume en Twitter un fío
de frases célebres dun tal David Castaño que me fixo moita graza. Daquela,
evidentemente, de ningún dos docentes, non tiña nin idea de quen eran, aínda
que si que me serviu para facerme unha idea do que me estaba esperar nos
vindeiros anos.

O conto das frases quedou aí, e non volveu aparecer, ata que comezou o 2º
cuadrimestre de primeiro. Nesa primeira semana, a comezos de febreiro de 2022,
coñecín a un profesor moi simpático chamado Alfredo Amigo\footnote{Alfredo, se
estás lendo isto tomando unhas garimbas na cafetaría, que che aproveiten!},
entre outros. O sábado seguinte, 5 de febreiro, era unha tarde fría de inverno
e dixen, vou apuntar frases. E aí comezou todo. Durante ese mesmo cuadrimestre
fun apuntando máis frases de profesores, onde destacaban como supoñeredes, as
de Eva Acosta, se cadra publicando algunha parvada de máis, coa inmadurez
lóxica dun rapaz de 18 anos.

Non obstante, se es Eva Acosta, Toni ou Alfredo Amigo e estás lendo isto,
lamento dicirche que non vas poder ler as túas propias frases. A comezos de
setembro de 2022, cando iniciei sesión en Twitter no portátil que teño cando
estou en Compostela, deume erro, e non puiden acceder máis á miña propia conta,
tiven que crear unha nova. Por tanto, todas as frases que apuntei durante o
primeiro curso non as ides poder ler, xa que non se encontran dispoñibles.

No entanto, o resto das xoias que hai, isto é, as de segundo, terceiro e
cuarto, si que están abertas ao público. Farei aquí un resumo a modo de
recompilación, incluíndo tamén algunha frase de primeiro que fun quen de
recuperar:

\Cita[Carlos Merino,\\ Física de Partículas Elementais, 2025.]{Todos tivemos
unha crisis cando fomos adolescentes e por iso decidimos estudar Física.}

\Cita[Carlos Merino,\\ Física de Partículas Elementais, 2025.]{O tempo é a
inversa da enerxía. Canto máis vello te fas, menos enerxía tes.}

\Cita[Carlos Merino,\\ Física de Partículas Elementais, 2025.]{Claro, é que
esta frase non está moi ben redactada e eso que a redactei eu, eh!?
Unbelievable.}

\Cita[Santos,\\ Mecánica Clásica III, 2023.]{De Métodos V sabedes un huevo de
curvas e superficies. Moito máis ca min, porque esa parte eu non a din.}

\Cita[Eva Acosta,\\ Métodos Matemáticos IV, 2022.]{Dos por uno, uno.
¡Jajajajaja! ¡Perdón! ¡Dos por uno, dos! ¡Yo no he dicho dos por uno, uno! ¡Lo
negaré!}

\Cita[Eva Acosta,\\ Métodos Matemáticos IV, 2022.]{Al de Métodos lo llamábamos
épsilon porque era pequeñito y despreciable.}

\Cita[Enrique Zas,\\ Física Xeral II, 2022.]{Vamos a suponer que esta pared es
rígida, lo cual no es cierto del todo, puesto que no es la primera vez que le
doy un golpe y vibra.}

\Cita[Jaime Álvarez,\\ Mecánica Clásica I, 2022.]{La energía mundial ni se crea
ni se destruye, solo se encarece. Esta es la nueva Ley de Conservación de la
Energía Mundial.}

\Cita[Jaime Álvarez,\\ Mecánica Clásica I, 2022.]{Mi coeficiente de reflexión
está tendiendo a cero. Cada vez reflexiono menos, pero cada vez transmito más.}

\Cita[Jaime Álvarez,\\ Mecánica Clásica I, 2022.]{Todos los días por estas
fechas hay estudiantes observando charcos. Esta es una buena frase para
Twitter.}

\Cita[Paco Ares,\\ Electromagnetismo II, 2023.]{Un teorema matemático no
depende de las condiciones en las que se está demostrando, a no ser que el
matemático esté fumado o drogado.}

%\Cita{El que respire se va a la calle.}{Curra,\\ Termodinámica e Teoría Cinética,
%2023.}

%O contexto desta última frase era que iamos facer un deses
%mitiquísimos controis sorpresa con Curra, que facía durante
%os últimos vinte minutos de clase. Só os máis veteranos desta
%facultade podemos dicir que tivemos o pracer (ou a desgraza)
%de termos sido alumnos desta profesora.

\Cita[Bernardo Adeva,\\ Física Cuántica I, 2023.]{Si los fotones estuvieran
sometidos al Principio de Exclusión de Pauli no se podría derretir un tanque
con un láser.}

% \Cita{Tiras un misil y que caiga con una precisión de 10 cm, por ejemplo, en el
% despacho de Vladímir Putin. [...] Sí, ha matado a 300 000 rusos, muchos
% jóvenes como vosotros, así que no vendría nada mal...}{Bernardo Adeva,\\ Física
% Cuántica I, 2023.}

\Cita[Suso Liñares,\\ Óptica II, 2024.]{La delta de Dirac se puede decir que
funciona de esa forma, aunque es una aberración. Pero vamos a usarla de esa
forma.}

\Cita[Manuel Caamaño,\\ Física Nuclear e de Partículas, 2024.]{Marie Curie es
la única persona con dos premios Nobel diferentes: Química y Física. Digo
premios Nobel de verdad, porque hay gente con premio Nobel de Química y de la
Paz y cosas así.}

\Cita[Manuel Caamaño,\\ Física Nuclear, 2025.]{Cuando fui a ver Oppenheimer al
cine estaba mi novia diciendo que me callara todo el tiempo. Yo estaba
<<¡BUAAAA, ESO TAMBIÉN LO VIMOS EN CLASE!>>.}

\Cita[Jaime Álvarez,\\ Astrofísica e cosmoloxía, 2025.]{Hay dos opciones:
aceptamos el principio cosmológico y seguimos dando clase o no lo aceptamos y
marchamos. Si queréis no lo aceptamos y nos vamos.}

% \Cita{Resulta que lleva miles de millones de años viajando y el pobre fotón
% acaba en uno de los detectores de aquí, del laboratorio de Óptica.}{(sen autoría aparente)}

\Cita[Jaime Álvarez,\\ Astrofísica e Cosmoloxía, 2025.]{Podía acabar en un
sitio mejor. Como yo, yo también podría acabar en un sitio mejor.}

%\Cita{Tienes que sumar 4+2=6, tienes que hacer bien la suma...}{Bernardo Adeva,\\
%Física de Partículas Elementais, 2025.}


\end{multicols}
